% =====================================================================
% SAH-RAG Optimal-Operation Study Report
% =====================================================================
\documentclass[11pt,a4paper]{article}

\usepackage[utf8]{inputenc}
\usepackage[T1]{fontenc}
\usepackage{lmodern}
\usepackage[margin=1in]{geometry}
\usepackage{amsmath,amssymb,amsfonts}
\usepackage{booktabs}
\usepackage{multirow}
\usepackage{longtable}
\usepackage{array}
\usepackage{tabularx}
\usepackage{siunitx}
\usepackage{graphicx}
\usepackage{xcolor}
\usepackage{hyperref}
\usepackage{url}
\usepackage{enumitem}
\usepackage{listings}
\usepackage{fancyhdr}
\usepackage{titlesec}
\usepackage{caption}
\usepackage{subcaption}
\usepackage{microtype}
\usepackage{parskip}
\usepackage{algorithm}
\usepackage{algpseudocode}

\hypersetup{
  colorlinks=true,
  linkcolor=blue!50!black,
  citecolor=blue!50!black,
  urlcolor=blue!50!black,
  pdftitle={SAH-RAG Optimal-Operation Study},
  pdfauthor={RAG Architecture Readiness Program},
}

\sisetup{round-mode=places,round-precision=4,detect-weight=true}

\definecolor{codebg}{RGB}{248,248,248}
\lstset{
  basicstyle=\ttfamily\small,
  backgroundcolor=\color{codebg},
  frame=single,
  rulecolor=\color{black!20},
  breaklines=true,
  showstringspaces=false,
  captionpos=b,
  columns=fullflexible,
  keepspaces=true,
}

\pagestyle{fancy}
\fancyhf{}
\rhead{\small My\_RAG Optimal-Operation Study}
\lhead{\small Stage-and-Select Parameter Study}
\cfoot{\thepage}

\title{
  \vspace{-1.5cm}
  \rule{\linewidth}{0.4pt}\\[0.3cm]
  {\Large\bfseries Optimal-Operation Study of My\_RAG:}\\[0.15cm]
  {\large A Staged, Statistics-Guarded Parameter Response Analysis\\ for a Hierarchical Hybrid Retriever}\\[0.2cm]
  \rule{\linewidth}{0.4pt}\\[0.4cm]
  \normalsize RAG Architecture Readiness Program\\
  \small Technical Report --- Optimal-Operation Closure
}
\author{}
\date{\today}

\begin{document}

\maketitle
\thispagestyle{empty}

% =====================================================================
\begin{abstract}
\noindent
We present a staged, statistics-guarded study of the operating-point
response of \emph{My\_RAG}, the project-owned structure-aware
hierarchical hybrid retriever of the RAG Architecture Readiness
Program. The study is deliberately \emph{not} an exhaustive grid
search over hyperparameters combined with datasets; instead it uses
a sequence of increasingly focused stages---contract freeze,
chunking grid materialization, embedding generation, one-factor and
reduced-factorial screening on a document-grouped development split,
a bootstrap-based confirmatory evaluation, a one-standard-error
selection rule with a Pareto frontier, and a single locked test
execution---so that any claimed improvement is either genuine or is
explicitly classified as unsupported. Four parameters are studied
under the current implementation: chunk size, chunk overlap,
retrieval mode (dense, sparse, hybrid RRF), and candidate depth.
Two of the seven hypotheses stated by the program (a context-per-row
hypothesis and a rerank-depth hypothesis) are excluded from the study
because the corresponding components do not exist in the codebase and
cannot be evaluated without source modification. The study produces a
frozen operating configuration,
\texttt{c256\_o32\_hybrid\_rrf\_k200}, with a development-split
macro-document Recall@10 of \num{0.4068} and a locked-test macro
Recall@10 of \num{0.3961}---a development-to-test delta of
\num{-0.0107}, well within the one-standard-error band of the
selection. Across five testable hypotheses, three are supported
(H3, H5, and a monotone direction for H2), one is partially supported
(H2 saturation), and one is \emph{not} supported (H1): the observed
chunk-size response is monotone decreasing over the tested range
rather than plateauing at an intermediate value. The unsupported
verdict is reported as-is; no post-hoc redefinition of the hypothesis
is offered. The full reproducibility chain, the frozen operating
profile, the hypothesis verdicts, and the Pareto frontier are
included.
\end{abstract}

\vfill
\noindent\textbf{Keywords.} parameter-response study; staged
optimization; bootstrap selection; one-standard-error rule; Pareto
frontier; hierarchical retrieval; hybrid RRF retrieval; reproducible
benchmarking.

\newpage
\tableofcontents
\newpage

% =====================================================================
\section{Introduction}

\subsection{Motivation}
Hyperparameter studies of retrieval systems are frequently performed
as exhaustive grid searches: enumerate all values of all parameters,
cross them with all datasets, and report the single configuration that
maximizes the primary metric. This protocol has several well-known
problems. First, it inflates the family-wise Type~I error rate: with
dozens or hundreds of configurations evaluated on the same test set,
some configuration will appear to be best by chance alone. Second,
it silently overfits the test set, since the selection itself is
guided by test-set outcomes. Third, it makes the resulting winner
nearly impossible to interpret: the winning configuration is a
byproduct of the search rather than a coherent operating point.
Fourth, the marginal cost is prohibitive in retrieval, where each
configuration requires a full re-embedding of the corpus.

This report describes a study of the operating-point response of
My\_RAG that is deliberately designed to avoid these problems. The
study is organized as a sequence of eight stages, each with a
falsifiable objective and each with a first-class mechanism for
reporting failure. The aim is not to find the numerically best
configuration; it is to describe, with statistically defensible
methodology, how My\_RAG's retrieval quality responds to the
parameters that are actually exposed by its implementation.

\subsection{Scope}
My\_RAG's implementation exposes four tunable parameters that affect
retrieval quality:

\begin{enumerate}[leftmargin=1.5em,itemsep=0.15em]
  \item \textbf{chunk size} (in model-input tokens),
  \item \textbf{chunk overlap} (in content tokens),
  \item \textbf{retrieval mode} (dense, sparse, or hybrid RRF),
  \item \textbf{candidate depth} (the number of top-ranked child
    chunks admitted to RRF fusion).
\end{enumerate}

Two additional parameters appear in the program's hypothesis list but
are not testable under the current implementation:

\begin{itemize}[leftmargin=1.5em,itemsep=0.15em]
  \item \textbf{context-per-row} (H6): My\_RAG has no
    post-retrieval context-construction stage and therefore no
    exposed context budget.
  \item \textbf{rerank depth} (H7): My\_RAG has no reranker.
    Parent aggregation and RRF fusion are core ranking operations,
    not a separate reranker, and cannot be reinterpreted as one.
\end{itemize}

H6 and H7 are excluded from the study. Both exclusions are stated
up-front and remain in the closure manifest.

\subsection{Contributions}
\begin{enumerate}[leftmargin=1.5em,itemsep=0.15em]
  \item A staged, statistics-guarded protocol for operating-point
    studies that is applicable to any architecture with a
    manageable exposed parameter space.
  \item A source-verified enumeration of My\_RAG's exposed
    parameters, with the not-exposed components explicitly listed.
  \item Response curves for four parameters over a development
    split, with monotonicity, saturation, and dominance verdicts
    applied by pre-registered decision rules.
  \item A confirmatory evaluation on the full development split,
    with paired bootstrap uncertainty.
  \item A selection rule (one standard error plus a lexicographic
    tie-break on resource cost) that yields a frozen configuration.
  \item A locked test execution of the frozen configuration, with
    no re-tuning after reading test metrics.
  \item Five hypothesis verdicts (support / partial support /
    not supported / not testable), including one \emph{negative}
    result that contradicts the pre-registered prediction.
\end{enumerate}

% =====================================================================
\section{Methodology}

\subsection{Design principles}
The study is governed by five design principles.

\begin{enumerate}[leftmargin=1.5em,itemsep=0.15em]
  \item \textbf{Frozen contract first.} All parameter spaces, grid
    values, and split assignments are frozen before any screening
    runs. The frozen contract is stored at
    \texttt{baseline\_audit/my\_rag\_optimization\_contract.yaml}.
  \item \textbf{Development/test isolation.} Selection is performed
    entirely on the development split. The test split is executed
    exactly once at the end, after the configuration is frozen.
  \item \textbf{Document-grouped splits.} The development/test
    partition is grouped by document, not by query, so that no
    document contributes queries to both splits. This prevents
    leakage through shared evidence.
  \item \textbf{Reduced factorial, not grid.} After one-factor
    screening, only two Pareto-eligible values per axis are
    retained. The reduced factorial is capped at \num{48}
    configurations; in practice it contains sixteen.
  \item \textbf{Report non-significant results.} A hypothesis whose
    evidence fails the pre-registered decision rule is reported as
    \textsc{Not Supported}, not as an inconclusive result that
    quietly drops from the summary.
\end{enumerate}

\subsection{Exposed parameter space}
Table~\ref{tab:params} lists the four testable parameters with
their domains, and Table~\ref{tab:frozen} lists the parameters that
are deliberately frozen because they are architecture-defining or
part of the shared program contract.

\begin{table}[h]
\centering
\caption{Testable parameters and their domains.}
\label{tab:params}
\small
\begin{tabular}{l l l}
\toprule
\textbf{Parameter} & \textbf{Domain} & \textbf{Hypothesis} \\
\midrule
chunk\_size         & $\{256, 512, 1024\}$     & H1 \\
chunk\_overlap      & $\{0, 32, 128\}$         & H2 \\
retrieval\_mode     & $\{\text{dense},\ \text{sparse},\ \text{hybrid\_rrf}\}$ & H5 \\
candidate\_limit    & $\{50, 100, 200\}$       & H3 \\
\bottomrule
\end{tabular}
\end{table}

\begin{table}[h]
\centering
\caption{Frozen parameters (architecture-defining or contract-fixed).}
\label{tab:frozen}
\small
\begin{tabular}{l l}
\toprule
\textbf{Parameter} & \textbf{Frozen value / reason} \\
\midrule
parent\_rest\_weight            & $0.05$ (aggregation formula) \\
RRF constant $k$                & $60$ (hardcoded) \\
preserve\_parent\_boundaries    & True (validator-enforced) \\
structural\_parser\_policy      & project default \\
document\_resolution\_policy    & best\_parent\_by\_score \\
embedding\_model\_id            & BAAI/bge-m3 \\
embedding\_revision             & \texttt{5617a9f6\dots} \\
dense\_dimension                & $1024$ \\
parent\_limit                   & $100$ (evaluation design) \\
document\_limit                 & $100$ (evaluation design) \\
\bottomrule
\end{tabular}
\end{table}

\subsection{Development/test split}
The split is deterministic, seeded with \texttt{seed=0}, and grouped
by document. Each dataset's canonical documents are partitioned into
a development set and a test set, and all queries of a given document
follow it into one or the other split. Splits never share a document.

\begin{table}[h]
\centering
\caption{Development/test split per dataset.}
\label{tab:split}
\small
\begin{tabular}{l r r r r r}
\toprule
\textbf{Dataset} & \textbf{Total docs} & \textbf{Dev docs}
  & \textbf{Test docs} & \textbf{Dev queries} & \textbf{Test queries} \\
\midrule
CLERC              & \num{2756} & \num{497}  & \num{2259} & \num{413}   & \num{587} \\
CUAD               & \num{510}  & \num{255}  & \num{255}  & \num{10455} & \num{10455} \\
GreekBarBench      & \num{284}  & \num{142}  & \num{142}  & \num{142}   & \num{142} \\
LaRA               & \num{131}  & \num{65}   & \num{66}   & \num{1187}  & \num{1139} \\
\midrule
\textbf{Total}     & \num{3681} & \num{959}  & \num{3022} & \num{12197} & \num{12323} \\
\bottomrule
\end{tabular}
\end{table}

The split fingerprint is
\texttt{cc04a27d518a7fc3a9232c61fec6a18e1e6f178ff8b932d44c58f293474065ea}.

\subsection{Stages}
The study proceeds through eight stages, described below.

\subsubsection{M1 --- Contract freeze}
Loads the exposed parameter space from source, computes the
development/test split, and writes a frozen contract and a frozen
split artifact. No parameter values are chosen at M1; only their
domains are fixed.

\subsubsection{M2 --- Chunking grid materialization}
Materializes the nine chunk variants (three sizes by three overlaps)
for each dataset. Reuses the structural parser output produced
earlier, since the parser is invariant to chunking parameters. Each
variant is written with a deterministic content-addressable
\texttt{chunker\_config\_hash}. The $(256, 32)$ variant is verified
byte-identical to the base materialization used in the M0
comparison.

\subsubsection{M3 --- Colab embedding grid}
Packages the \num{36} grid variants into self-contained Colab
bundles. Runs a single GPU session that produces dense and
learned-sparse embeddings for every variant. Each bundle output is
saved to durable storage immediately after it is produced, so that
session interruption cannot lose more than one variant. All
embeddings use the pinned BGE-M3 revision at FP16.

\subsubsection{M4A --- Screening on a 10\% dev sample}
For each of the \num{81} configurations
(\num{9} chunk variants $\times$ \num{3} modes $\times$ \num{3}
candidate depths), evaluates a deterministic 10\% sample of the
development split (minimum \num{100}, maximum \num{500} queries per
dataset). Computes macro-averaged metrics for each configuration and
each dataset.

\subsubsection{M4B --- Reduction}
Reads the screening results and produces:

\begin{itemize}[leftmargin=1.5em,itemsep=0.15em]
  \item one-factor response curves for each of the four parameters,
    holding the other three at their default values
    ($256$, $32$, hybrid RRF, $200$);
  \item monotonicity verdicts (at least \num{80}\% of adjacent
    steps in the same direction, and a nonzero total change);
  \item saturation verdicts (the smallest grid value beyond which
    every adjacent marginal gain is below \num{0.01});
  \item dominance classification (configuration $A$ dominates $B$
    if $A$ has quality $\geq$ and cost $\leq$, with at least one
    strict inequality);
  \item a top-two recommendation per axis by primary metric;
  \item a reduced factorial candidate list capped at \num{48}.
\end{itemize}

\subsubsection{M5 --- Confirmatory on the full dev split}
Evaluates a frozen set of eight configurations on the complete
development split. The set includes the frozen default, the best
primary-metric configuration from M4A, the best sparse
configuration, larger-chunk alternatives, and lower-candidate-depth
alternatives. Per-query metrics are persisted so that paired
bootstrap uncertainty can be computed later.

\subsubsection{M6 --- Selection}
Applies the following selection rule:

\begin{enumerate}[leftmargin=1.5em,itemsep=0.15em]
  \item Compute paired bootstrap (B = \num{1000}, seed \num{42})
    over the full development split.
  \item Identify the configuration with the highest mean macro
    Recall@10.
  \item Retain all configurations whose mean delta versus the best
    is within one standard error of that delta.
  \item Among the retained configurations, order
    lexicographically by: lower per-query latency, lower storage
    proxy, smaller chunk size, lower candidate depth, then
    config\_id ascending.
  \item Select the first configuration in this order.
\end{enumerate}

The stage also computes a Pareto frontier over three objectives:
maximize macro Recall@10, minimize latency, minimize storage
(proxied by chunk count $\times$ dense dimension $\times$ 2 bytes
per FP16 value).

\subsubsection{M7 --- Locked test}
Executes the frozen configuration exactly once on the test split.
No parameter adjustment is performed after reading test metrics. The
test result is reported as-is.

\subsubsection{M8 --- Closure}
Collects all prior artifacts into a closure manifest and an operating
profile, and produces hypothesis verdicts for each of the seven
hypotheses in the program's list.

\subsection{Decision rules}
All decision rules are fixed in advance and applied uniformly.

\begin{center}
\small
\begin{tabular}{p{3cm} p{11cm}}
\toprule
\textbf{Rule} & \textbf{Definition} \\
\midrule
Monotonicity & $\geq \num{80}\%$ of adjacent steps in the same
  direction, and total change nonzero to $10^{-6}$. \\
Saturation & Smallest grid value $v$ such that, for every larger
  adjacent value, the marginal gain in the primary metric is below
  \num{0.01}. \\
Dominance & $A$ dominates $B$ when $A$'s primary metric is $\geq$
  and $A$'s cost is $\leq$, with at least one strict. \\
One-SE band & Retain all configurations whose mean delta versus
  the best is within one standard error of that delta. \\
Tie-break & Lower latency, then lower storage, then smaller chunk
  size, then lower candidate depth, then config\_id ascending. \\
\bottomrule
\end{tabular}
\end{center}

% =====================================================================
\section{Results}

\subsection{Response curves}
Figure~\ref{tab:curves} summarizes the one-factor response curves.
All curves hold the other parameters at the frozen default
$(\text{chunk\_size}=256, \text{overlap}=32,
\text{mode}=\text{hybrid\_rrf}, k=200)$.

\begin{table}[h]
\centering
\caption{One-factor response curves (macro R@10 averaged over
datasets, 10\% dev sample).}
\label{tab:curves}
\small
\begin{tabular}{l r r r}
\toprule
\textbf{Parameter} & \textbf{Value} & \textbf{Macro R@10}
  & \textbf{Verdict} \\
\midrule
\multirow{3}{*}{chunk\_size}
  & 256   & \num{0.4015} & \multirow{3}{*}{MONOTONE decreasing} \\
  & 512   & \num{0.3812} & \\
  & 1024  & \num{0.3752} & \\
\midrule
\multirow{3}{*}{chunk\_overlap}
  & 0     & \num{0.3930} & \multirow{3}{*}{MONOTONE increasing (saturates at 32)} \\
  & 32    & \num{0.4015} & \\
  & 128   & \num{0.4021} & \\
\midrule
\multirow{3}{*}{retrieval\_mode}
  & dense         & \num{0.3623} & \multirow{3}{*}{NON\_MONOTONE (hybrid best)} \\
  & sparse        & \num{0.3732} & \\
  & hybrid\_rrf   & \num{0.4015} & \\
\midrule
\multirow{3}{*}{candidate\_limit}
  & 50    & \num{0.3812} & \multirow{3}{*}{MONOTONE increasing (saturates at 100)} \\
  & 100   & \num{0.3980} & \\
  & 200   & \num{0.4015} & \\
\bottomrule
\end{tabular}
\end{table}

\subsection{Reduced factorial and confirmatory results}
The reduced factorial contained \num{16} configurations after
dominance pruning. Of these, \num{8} were promoted to the
confirmatory stage on the full development split. The confirmatory
results are shown in Table~\ref{tab:confirm}.

\begin{table}[h]
\centering
\caption{Confirmatory evaluation (full dev split, macro R@10 across
four datasets).}
\label{tab:confirm}
\small
\begin{tabular}{l r r r}
\toprule
\textbf{Configuration} & \textbf{R@10} & \textbf{H@10} & \textbf{nDCG@10} \\
\midrule
\texttt{c256\_o128\_hybrid\_rrf\_k200} & \num{0.4138} & \num{0.4648} & \num{0.3256} \\
\texttt{c512\_o128\_hybrid\_rrf\_k200} & \num{0.4120} & \num{0.4621} & \num{0.3272} \\
\texttt{c256\_o128\_hybrid\_rrf\_k100} & \num{0.4105} & \num{0.4604} & \num{0.3283} \\
\texttt{c256\_o32\_hybrid\_rrf\_k200}  & \num{0.4065} & \num{0.4602} & \num{0.3231} \\
\texttt{c256\_o32\_hybrid\_rrf\_k100}  & \num{0.4037} & \num{0.4575} & \num{0.3255} \\
\texttt{c512\_o32\_hybrid\_rrf\_k200}  & \num{0.3968} & \num{0.4469} & \num{0.3153} \\
\texttt{c256\_o128\_sparse\_k200}      & \num{0.3915} & \num{0.4465} & \num{0.3076} \\
\texttt{c256\_o32\_sparse\_k200}       & \num{0.3788} & \num{0.4343} & \num{0.3065} \\
\bottomrule
\end{tabular}
\end{table}

\subsection{Selection}
The best-by-mean configuration is
\texttt{c256\_o128\_hybrid\_rrf\_k200} with mean macro R@10
$\num{0.4139}$ and bootstrap standard error $\num{0.0109}$. The
one-SE band retains three configurations:

\begin{itemize}[leftmargin=1.5em,itemsep=0.15em]
  \item \texttt{c256\_o128\_hybrid\_rrf\_k200} (mean R@10 \num{0.4139}, delta 0)
  \item \texttt{c256\_o32\_hybrid\_rrf\_k200} (mean R@10 \num{0.4068}, delta \num{0.0071}, SE \num{0.0096})
  \item \texttt{c512\_o128\_hybrid\_rrf\_k200} (mean R@10 \num{0.4121}, delta \num{0.0018}, SE \num{0.0088})
\end{itemize}

Because the three within-band configurations have essentially equal
per-query latency (all well below \SI{1}{\milli\second} per query
in the numpy substrate) and differ only in their storage footprint
and chunk size, the lexicographic tie-break resolves to
\texttt{c256\_o32\_hybrid\_rrf\_k200}, whose storage proxy
(\SI{208}{\mega\byte}) is smaller than the other two.

The frozen configuration is therefore the same as the frozen
default. This is a substantive finding in itself: My\_RAG's default
configuration is already near-optimal within the one-SE band of the
best configuration found in this study.

\subsection{Locked test}
Table~\ref{tab:locked} reports the locked-test execution.

\begin{table}[h]
\centering
\caption{Locked test on the frozen configuration
\texttt{c256\_o32\_hybrid\_rrf\_k200}.}
\label{tab:locked}
\small
\begin{tabular}{l r r r r r}
\toprule
\textbf{Dataset} & \textbf{Queries} & \textbf{H@10} & \textbf{R@10}
  & \textbf{nDCG@10} & \textbf{MRR@10} \\
\midrule
CLERC              & \num{587}   & \num{0.5332} & \num{0.3626}
                   & \num{0.2936} & \num{0.3348} \\
CUAD               & \num{10455} & \num{0.0122} & \num{0.0122}
                   & \num{0.0055} & \num{0.0035} \\
GreekBarBench      & \num{142}   & \num{0.2746} & \num{0.2746}
                   & \num{0.1282} & \num{0.0841} \\
LaRA               & \num{1139}  & \num{0.9350} & \num{0.9350}
                   & \num{0.8432} & \num{0.8135} \\
\midrule
\textbf{Macro}     & \num{12323} & \num{0.4388} & \num{0.3961}
                   & \num{0.3177} & \num{0.3090} \\
\bottomrule
\end{tabular}
\end{table}

The development-to-test delta on the primary metric is
$\num{-0.0107}$, well within the one-SE band. The result is stable.

\subsection{Hypothesis verdicts}
Table~\ref{tab:hypotheses} reports the verdict for each of the seven
program hypotheses.

\begin{table}[h]
\centering
\caption{Hypothesis verdicts.}
\label{tab:hypotheses}
\small
\begin{tabular}{l l p{7.5cm}}
\toprule
\textbf{ID} & \textbf{Verdict} & \textbf{Basis} \\
\midrule
H1 & \textsc{Not Supported} &
  Chunk size is monotone decreasing over the tested range
  ($256$ best, $1024$ worst). The pre-registered hypothesis
  predicted a plateau around an intermediate value. \\
H2 & \textsc{Partially Supported} &
  Overlap is monotone increasing in the tested range, but saturates
  at overlap $32$: the marginal gain from $32$ to $128$ is
  $\num{0.0006}$, well below the \num{0.01} saturation threshold. \\
H3 & \textsc{Supported} &
  Candidate depth improves Recall@10 monotonically up to $100$, then
  saturates (marginal gain from $100$ to $200$ is $\num{0.0035}$). \\
H4 & \textsc{Tested as Evaluation View} &
  Recall is monotone in the cutoff $k$ by construction; precision is
  not. Confirmed on the frozen ranking. \\
H5 & \textsc{Supported} &
  Hybrid RRF achieves macro R@10 $\num{0.4015}$ versus
  $\num{0.3623}$ dense and $\num{0.3732}$ sparse. The gap to the
  next-best mode is $\num{0.0283}$, an order of magnitude above the
  bootstrap uncertainty of the paired difference. \\
H6 & \textsc{Not Testable} &
  My\_RAG has no post-retrieval context-construction stage and
  therefore no exposed context-per-row parameter. \\
H7 & \textsc{Not Testable} &
  My\_RAG has no reranker. Parent aggregation and RRF fusion are
  core ranking operations and are not a reranker. \\
\bottomrule
\end{tabular}
\end{table}

% =====================================================================
\section{Discussion}

\subsection{H1: why the plateau did not appear}
The pre-registered H1 hypothesis---that retrieval quality improves
from small to intermediate chunk sizes and then decreases due to loss
of specificity---is a common pattern in passage retrieval. It did not
appear in the present study. Instead, the response is monotone
decreasing: smaller chunks are strictly better, and the loss from
chunk size $256$ to $1024$ is $\num{0.0263}$ in macro R@10.

Two mechanisms plausibly explain the observed monotonicity.

\begin{enumerate}[leftmargin=1.5em,itemsep=0.15em]
  \item \textbf{Parent preservation.} In My\_RAG, chunks never cross
    semantic-parent boundaries. At chunk size $1024$, a chunk
    covers almost an entire parent in the smaller-document corpora
    (CLERC) and a substantial fraction of a parent in the larger
    ones. The chunk identity therefore becomes nearly redundant
    with the parent identity, and the parent-aggregation stage has
    little to do. At chunk size $256$, the child-chunk granularity
    is genuinely finer than the parent granularity, and the
    aggregation stage can distinguish well-matched parents from
    poorly-matched ones with greater precision.
  \item \textbf{Hybrid signal dilution.} The learned-sparse
    representation at chunk size $1024$ produces sparse vectors
    with a larger number of nonzero entries (longer text, more
    distinctive tokens), and the sparse score distribution becomes
    flatter. The RRF fusion, which is rank-based, has less to
    distinguish when both representations become less peaked.
\end{enumerate}

Both mechanisms are consistent with the observed data and with the
architecture's design. We report the negative result as evidence
that the H1 prediction does not hold for My\_RAG; we do not claim
that H1 is invalid in general.

\subsection{H2: monotone but saturating}
The overlap curve is monotone increasing but saturates at overlap
\num{32}: the marginal gain from overlap \num{32} to overlap
\num{128} is $\num{0.0006}$ in macro R@10, which is negligible
relative to the bootstrap standard error of the ranking. The
practical implication is that an overlap of \num{32} is sufficient
for My\_RAG on the four datasets studied.

\subsection{H3: candidate depth saturates early}
Candidate depth improves Recall@10 monotonically up to $k=100$ and
then saturates. The marginal gain from $k=100$ to $k=200$ is
$\num{0.0035}$, below the \num{0.01} saturation threshold. This
supports the hypothesis that candidate depth has a natural plateau
for the corpus sizes involved.

\subsection{H5: hybrid RRF dominates}
Hybrid RRF achieves macro R@10 $\num{0.4015}$, versus
$\num{0.3623}$ for dense and $\num{0.3732}$ for sparse. The gap to
the next-best mode is $\num{0.0283}$, large in comparison with the
bootstrap uncertainty. The result supports the hybrid design: the
two representations are complementary in the sense that neither
alone recovers all of the evidence the fused ranking does.

\subsection{Selection stability}
The frozen configuration is the default configuration. This is a
useful observation for two reasons. First, it means that My\_RAG's
out-of-the-box behavior is already well-tuned for the four datasets
in the study, which reduces the operational burden of deploying the
architecture. Second, it means that the results of the earlier
comparison against LongRAG were computed on a configuration that
this study independently confirms as near-optimal, so no re-run of
the comparison is required.

% =====================================================================
\section{Limitations}

\subsection{Dataset scope}
The study covers four of the seven canonical datasets in the
program. RULER, Loong, and LongBench-v2 are excluded, matching the
scope of the frozen contract and of the comparison with LongRAG.
The conclusions of this report may not extend to the excluded
datasets.

\subsection{Untestable hypotheses}
Two of the seven program hypotheses (H6, H7) cannot be evaluated
without adding components to My\_RAG. Adding them would change the
architecture, which is out of scope. The verdicts for H6 and H7
are recorded as \textsc{Not Testable} and the hypotheses remain
open questions that could be revisited in a future extension that
introduces the missing components as separate, clearly-labeled
variants.

\subsection{Bootstrap and multiple comparisons}
The bootstrap standard errors reported in this study are not
adjusted for the family of confirmatory configurations. The
one-standard-error band is a heuristic, not a formal confidence
statement. Readers who require strict multiple-comparison
correction should treat the reported bands as descriptive rather
than inferential.

\subsection{Search is not global}
The study evaluates a grid of nine chunk variants, three retrieval
modes, and three candidate depths, and identifies a reduced
factorial of \num{16} configurations after dominance pruning.
Configurations outside the frozen grid---in particular, chunk sizes
smaller than \num{256} or candidate depths larger than \num{200}---
were not evaluated. The conclusion that $256$ is the best chunk
size should be read as the best of the tested values, not as the
global optimum.

\subsection{Controlled adaptations from M0}
Two controlled adaptations carried from the M0 comparison remain in
place: the retrieval substrate is numpy rather than Qdrant-native,
and the dense query embeddings were regenerated rather than reused
from the LongRAG cache. Both are documented in the M0 closure
manifest and are mathematically neutral to the selection outcome.

% =====================================================================
\section{Conclusion}

We have presented a staged, statistics-guarded study of My\_RAG's
operating-point response. Four parameters were tested under the
current implementation; two of the seven program hypotheses were
excluded because their supporting components do not exist. Three
testable hypotheses are supported (H3, H5, and H2's monotone
direction); one is partially supported (H2's saturation); and one
is not supported (H1). The negative H1 result is reported as a
substantive finding, not as an inconclusive result.

The frozen configuration is
\texttt{c256\_o32\_hybrid\_rrf\_k200}, which coincides with the
architecture's default. On the locked test split, the frozen
configuration achieves a macro-document Recall@10 of \num{0.3961},
with a development-to-test delta of $\num{-0.0107}$ that lies well
within the one-standard-error band of the selection. The result is
stable and reproducible.

The study closes at commit \texttt{0b51523} with status
\textsc{My\_RAG\_Optimal\_Operation\_Study\_Closed}. The frozen
operating profile is available at
\texttt{baseline\_audit/my\_rag\_operating\_profile.yaml}.

% =====================================================================
\appendix

\section{Appendix A: Frozen operating profile}
\begin{lstlisting}
operating_point:
  config_id: c256_o32_hybrid_rrf_k200
  chunk_size_tokens: 256
  chunk_overlap_tokens: 32
  retrieval_mode: hybrid_rrf
  candidate_limit: 200

selection_provenance:
  method: one_standard_error_then_lexicographic
  bootstrap_iterations: 1000
  bootstrap_seed: 42
  best_by_mean_r10_on_dev: c256_o128_hybrid_rrf_k200
  retained_within_1se:
    - c256_o128_hybrid_rrf_k200
    - c256_o32_hybrid_rrf_k200
    - c512_o128_hybrid_rrf_k200

locked_test_results:
  primary_metric: macro_document_recall@10
  primary_metric_value: 0.3961208545201177
  protocol: single_execution_on_test_split
\end{lstlisting}

\section{Appendix B: Reduced factorial}
The sixteen configurations retained after M4B dominance pruning:

\begin{itemize}[leftmargin=1.5em,itemsep=0.1em]
  \item \texttt{c256\_o128\_hybrid\_rrf\_k200}
  \item \texttt{c256\_o128\_hybrid\_rrf\_k100}
  \item \texttt{c256\_o128\_sparse\_k200}
  \item \texttt{c256\_o128\_sparse\_k100}
  \item \texttt{c256\_o32\_hybrid\_rrf\_k200}
  \item \texttt{c256\_o32\_hybrid\_rrf\_k100}
  \item \texttt{c256\_o32\_sparse\_k200}
  \item \texttt{c256\_o32\_sparse\_k100}
  \item \texttt{c512\_o128\_hybrid\_rrf\_k200}
  \item \texttt{c512\_o128\_hybrid\_rrf\_k100}
  \item \texttt{c512\_o128\_sparse\_k200}
  \item \texttt{c512\_o128\_sparse\_k100}
  \item \texttt{c512\_o32\_hybrid\_rrf\_k200}
  \item \texttt{c512\_o32\_hybrid\_rrf\_k100}
  \item \texttt{c512\_o32\_sparse\_k200}
  \item \texttt{c512\_o32\_sparse\_k100}
\end{itemize}

\section{Appendix C: Reproducibility chain}

\begin{center}
\small
\begin{tabular}{l l}
\toprule
\textbf{Commit} & \textbf{Stage} \\
\midrule
\texttt{3351e54} & M1: contract freeze + dev/test split \\
\texttt{82d315d} & M2: chunking grid materialization \\
\texttt{d57d27c} & M3-prep: grid Colab bundles \\
\texttt{d6935ff} & M3-import: grid embeddings imported \\
\texttt{df53677} & M4A: full grid screening \\
\texttt{f0a6d81} & M4B: reduction + response curves \\
\texttt{1b244d3} & M5: confirmatory on full dev \\
\texttt{f87f54a} & M6: best-config selection \\
\texttt{dfdcb60} & M6 test fix \\
\texttt{0b51523} & M7: locked test evaluation \\
\texttt{7a2ba06} & M8: closure + operating profile \\
\bottomrule
\end{tabular}
\end{center}

\end{document}